\subsection{Projective Modules over an Artinian Ring}

PROPOSITION 4. --- Every module over a left Artinian ring has a projective cover.

Let A be a left Artinian ring, and let M be a module over A. The radical r of A is a nilpotent two-sided ideal, and the ring A/r is semisimple (VIII, p.~173, Proposition 1). It follows that the A/r-module M/rM is projective. The assertion then follows from Proposition 11 of VIII, p.~164.

PROPOSITION 5. --- Let A be a left Artinian ring and r its radical.

\begin{enumerate}
\def\labelenumi{\alph{enumi})}
\item
  Let P be a projective A-module. Denote the canonical mapping from P to P/rP by u. Then \((P,u)\) is a projective cover of P/rP. In particular, the A-module P is finitely generated if and only if P/rP is.
\item
  Let M be a module over the ring A/r. Then M, viewed as an A-module, has a projective cover. If \((P,u)\) is such a projective cover, then when passing to the quotient, u induces an isomorphism from P/rP to M. Moreover, P is indecomposable if and only if M is simple.
\item
  Let M and M' be modules over the ring A/r. Let \((P, u)\) and \((P', u')\) be projective covers of the A-modules M and M'. Then M and M' are isomorphic if and only if P and P' are.
\end{enumerate}

The radical r of A is a nilpotent two-sided ideal (VIII, p.~173, Proposition 1). Assertion a) therefore follows from Corollary 1 of VIII, p.~163 and Remark 2 of VIII, p.~161.

Let us prove assertion b). The existence of a projective cover \((P, u)\) of M follows from Proposition 4 and the isomorphism from P/rP to M of Proposition 9 of VIII, p.~162. Since the ring A/r is semisimple, a module over this ring is indecomposable if and only if it is simple. The last assertion then follows from Corollary 2 of VIII, p.~160.

Let us prove c). Let f be an isomorphism from P to \(P'\) ; it induces an isomorphism \(\overline{f}\) from \(P/rP\) to \(P'/rP'\) . By b), these quotients are isomorphic to M and \(M'\) , respectively. Therefore, M is isomorphic to \(M'\) . Conversely, by Proposition 8 of VIII, p.~161, every isomorphism f from M to \(M'\) lifts to an isomorphism \(\widetilde{f}\) from P to \(P'\) such that \(f \circ u = u' \circ \widetilde{f}\) .

COROLLARY. --- With each class P of projective A-modules, associate the class \(\varphi(\mathrm{P})\) of the module P/rP over the ring A/r. Every class of modules (resp. of finitely generated modules) over the ring A/r is of the form \(\varphi(\mathrm{P})\) for a unique class P of projective A-modules (resp. of finitely generated projective A-modules).

Let A be a left Artinian ring and \(\mathfrak{r}\) its radical. Denote the set of classes of simple A-modules by \(\mathcal{S}\) ; it is finite (VIII, p.~51). By Proposition 2 of VIII, p.~174, it can be canonically identified with the set of classes of simple modules over the semisimple ring \(\mathrm{A} / \mathfrak{r}\) . For every \(\lambda \in \mathcal{S}\) , let \(\mathrm{S}_{\lambda}\) be a simple A-module of class \(\lambda\) , and choose a projective cover \((\mathrm{P}_{\lambda}, u_{\lambda})\) of \(\mathrm{S}_{\lambda}\) (VIII, p.~175, Proposition 4). By Proposition 5, the A-module \(\mathrm{P}_{\lambda}\) is projective and indecomposable, and \(u_{\lambda}\) defines an isomorphism from \(\mathrm{P}_{\lambda} / \mathfrak{r}\mathrm{P}_{\lambda}\) to \(\mathrm{S}_{\lambda}\) . Moreover, if P is a projective and indecomposable A-module, then there exists a unique \(\lambda \in \mathcal{S}\) such that P is isomorphic to \(\mathrm{P}_{\lambda}\) : it is the unique \(\lambda \in \mathcal{S}\) such that \(\mathrm{P} / \mathfrak{r}\mathrm{P}\) is isomorphic to \(\mathrm{S}_{\lambda}\) .

PROPOSITION 6. --- Let A be a left Artinian ring and r its radical. Let P be a projective A-module.

\begin{enumerate}
\def\labelenumi{\alph{enumi})}
\item
  The A-module \(\overline{\mathrm{P}} = \mathrm{P} / \mathfrak{r}\mathrm{P}\) is semisimple, and the A-module \(\mathrm{P}\) is isomorphic to \(\bigoplus_{\lambda \in \mathcal{S}}\mathrm{P}_{\lambda}^{([\overline{\mathrm{P}}:\lambda ])}\) .
\item
  Let \((\mathrm{Q}_{i})_{i\in\mathrm{I}}\) be a family of indecomposable projective submodules of P with direct sum P. Then for every \(\lambda\in S\) , the cardinal of the set \(\mathrm{I}(\lambda)\) of \(i\in I\) such that \(Q_{i}\) is isomorphic to \(P_{\lambda}\) is equal to \([\overline{P}:\lambda]\) .
\end{enumerate}

The fact that \(\overline{P}\) is semisimple follows from Proposition 2 (VIII, p.~174). The A-module \(Q = \oplus_{\lambda \in S} P_{\lambda}^{([\overline{P}: \lambda])}\) is projective. Since \(P_{\lambda}/rP_{\lambda}\) is isomorphic to \(S_{\lambda}\) , the quotient Q/rQ is isomorphic to \(\oplus_{\lambda \in S} S_{\lambda}^{([\overline{P}: \lambda])}\) , that is, to \(\overline{P} = P/rP\) . By Proposition 5, the A-modules P and Q are isomorphic.

Suppose given a family \((\mathrm{Q}_{i})_{i\in\mathrm{I}}\) of projective and indecomposable submodules with direct sum P. By Proposition 5, the A-module \(Q_{i}/rQ_{i}\) is simple for every \(i\in I\) . For \(\lambda\in S\) , denote by \(\mathrm{I}(\lambda)\) the set of \(i\in I\) such that \(Q_{i}/rQ_{i}\) is isomorphic to \(S_{\lambda}\) , hence to \(P_{\lambda}/rP_{\lambda}\) ; it is also the set of \(i\in I\) such that \(Q_{i}\) is isomorphic to \(P_{\lambda}\) . Since \(\overline{P}=P/rP\) is the direct sum of the family \((\mathrm{Q}_{i}/rQ_{i})_{i\in\mathrm{I}}\) , we have \(\operatorname{Card}(\mathrm{I}(\lambda))=[\overline{\mathrm{P}}:\lambda]\) by Theorem 1 of VIII, p.~32.

Example. --- Take P equal to \(A_{s}\) . For \(\lambda \in S\) , denote the opposite field of the commutant of the simple A-module \(S_{\lambda}\) by \(D_{\lambda}\) and the dimension of \(S_{\lambda}\) viewed as a right vector space over the field \(D_{\lambda}\) by \(m(\lambda)\) . We know that \(m(\lambda)\) is equal to the multiplicity \([A_{s}/rA_{s} : \lambda]\) (VIII, p.~143, Proposition 11). Consequently, the A-module \(A_{s}\) is isomorphic to \(\oplus_{\lambda \in S} P_{\lambda}^{m(\lambda)}\) .

